\documentclass[french]{book}
\usepackage{teach}
\usepackage{verbatim}
\usepackage{amsmath,amsfonts,amssymb}
\usepackage{pgfplots}
%\usepackage{geometry}
%\geometry{top=1cm, bottom=1cm, left=2cm, right=2cm}
\usepackage[utf8]{inputenc}
\usepackage[french]{babel}
\redefinecolor{thm}{title}{bgcolor}{alizarin}
\redefinecolor{prop}{title}{bgcolor}{meatbrown}
\redefinecolor{defn}{title}{bgcolor}{olivine}
\definecolor{alizarin}{rgb}{0.82, 0.1, 0.26}
\definecolor{meatbrown}{rgb}{0.9, 0.72, 0.23}
\definecolor{olivine}{rgb}{0.6, 0.73, 0.45}
\usepackage{pgf,tikz,tkz-tab}

\usepackage{mathrsfs}
\usetikzlibrary{arrows}

\begin{document}
\begin{center}
\huge{ Devoir surveillé VIII}
\end{center}
%\title{Devoir surveillé V}

\section*{Série statistique}

\begin{enumerate}
\item Donner une liste de 5 valeurs pour laquelle la médiane vaut 12 et la moyenne vaut 10.

\item Donner une liste de 10 valeurs pour laquelle la médiane est identique au premier quartile.
\item Donner deux séries distinctes comportant chacune 6 valeurs, ayant la même valeur médiane et les mêmes quartiles Q1 et Q3 mais dont l’une a 120 pour valeur maximale, alors que l’autre a 12 pour valeur maximale. \\
\end{enumerate}
 
\section*{Étude de salaire}

On considère la série suivante qui correspond à la répartition des salaires mensuels en euros dans une entreprise.

\begin{center}
   \begin{tabular}{| l || c | c | c | c | c |}
     \hline
     Salaires mensuels & 1 000 & 1 200 & 1 300 & 1 500 &  5000 \\ \hline
     Effectifs & 11 & 9 & 14 & 4 & 1 \\ 
     \hline
   \end{tabular}
 \end{center}

\begin{enumerate}
\item Calculer le salaire moyen et le salaire médian au sein de cette entreprise.

\item Déterminer les quartiles Q1 et Q3 en expliquant la démarche.
Interpréter concrètement la valeur trouvée pour Q3.

\item Pour vous recruter, le directeur de l’entreprise vous présente la valeur du salaire moyen. Est-ce un indicateur pertinent ici ? Expliquer.

\end{enumerate}


\section*{Extractions de données et interprétations}
Dans une maternité, on a relevé les masses \textit{(en kg)} à la naissance de 103 bébés nés en avril 2019.


\begin{center}
   \begin{tabular}{| l || c | c | c | c | c |}
     \hline
     Masse & 2,5 & 2,6 & 3 & 3,1 & 5 \\ \hline
     Effectif & 4 & 3 & 9 & 7 & 1 \\ \hline
	  Effectif cumulés & & & & & \\ \hline
	  Fréquence cumulées & & & & & \\
 
     \hline
   \end{tabular}
 \end{center}

\begin{enumerate}
\item Recopier le tableau et complèter les cases vides.
\item Par simple lecture ou calcul mental, combien de bébés pesaient moins de 3kg ? Plus de 3,1kg ? 
\item Quelle est la fréquence des bébés ne dépassant pas 3,1 kg ?
\item Quel est le pourcentage de bébés pesant plus de 2,5 kg et moins de 3,2 kg ? 
\end{enumerate}

\end{document}